\documentclass[11pt]{article}

\usepackage[margin=1in]{geometry}
\usepackage{amsmath,amssymb,amsthm,mathtools}
\usepackage{booktabs}
\usepackage{enumitem}
\usepackage[hidelinks]{hyperref}
\usepackage[nameinlink,capitalize]{cleveref}

\newtheorem{theorem}{Theorem}[section]
\newtheorem{lemma}[theorem]{Lemma}
\newtheorem{proposition}[theorem]{Proposition}
\newtheorem{corollary}[theorem]{Corollary}
\theoremstyle{definition}
\newtheorem{definition}[theorem]{Definition}
\newtheorem{remark}[theorem]{Remark}

\newcommand{\cube}[1]{\{0,1\}^{#1}}
\newcommand{\Ftwo}{\mathbb F_2}
\newcommand{\R}{\mathbb R}
\newcommand{\wt}{\operatorname{wt}}
\newcommand{\NAE}{\operatorname{NAE}}
\newcommand{\OR}{\operatorname{OR}}
\newcommand{\ANF}{\operatorname{ANF}}
\newcommand{\supp}{\operatorname{supp}}

\title{Exact Repair Obstructions for a Sensitivity--14 Composition\\
Built from the Kushilevitz Function}
\author{Jia Lei \and Bokun Zhao}
\date{August 14, 2026}

\begin{document}
\maketitle

\begin{abstract}
Eyal Kushilevitz's six-variable Boolean function, recorded in a footnote of
Nisan and Wigderson, has real degree three and full sensitivity six.  We use a
coordinate relabeling of that function as an outer map.  Substituting four
disjoint three-variable not-all-equal gates and two coordinate inputs gives an
explicit Boolean function on fourteen variables with sensitivity fourteen and
real degree exactly six.  Its degree-six homogeneous component is
\(-E_1E_3(E_0+E_2)\), consisting of exactly 54 monomials, each with
coefficient \(-1\).

The main results exclude several precisely specified routes to lowering this
construction to degree five.  First, a degree-five, sensitivity-twelve base
\(A\), built from the same outer function, admits no Boolean function \(B\) on
the same twelve variables satisfying \(B(0)=1\) and \(\deg(B-A)\le4\).  Second, a target-labelled
repair cannot factor only through the four not-all-equal outputs and the two
coordinate ports.  Third, twisting any such factored function by the XOR
parities of all four raw blocks always has degree greater than five.  Full
proofs rest on two direct-sum noncancellation lemmas: one for weighted top
signatures under disjoint block substitution, and one for blockwise Walsh
signatures.

We also give two exact, lossless parameterizations of all Boolean corrections
relative to the fixed base \(P\): an additive residual system and an XOR
correction system.  For the simplest XOR representative
\(J_0=g_1g_3(g_0\oplus g_2)\), one has \(\deg_{\mathbb R}J_0=8\), while exact
integer M\"obius inversion shows that \(P\oplus J_0\) has degree ten.  Its top
component is \(4E_0E_1E_2E_3x_{12}x_{13}\), comprising exactly 81 monomials.
A standalone verifier accompanies the paper.  The results do not construct a
degree-five, sensitivity-fourteen function and do not improve the Kushilevitz
recursive exponent.
\end{abstract}

\section{Introduction}

For a Boolean function \(f:\cube n\to\{0,1\}\), let
\(\deg_{\R}(f)\) be the degree of its unique real multilinear polynomial and
let \(s(f)\) be its maximum sensitivity.  Nisan and Szegedy established the
foundational connection between real degree and combinatorial query measures
\cite{NisanSzegedy1994}; Huang later proved the sensitivity conjecture
\cite{Huang2019}.  The explicit recursive benchmark used in this paper comes
from a six-variable function due to Eyal Kushilevitz.  Kushilevitz did not
publish the function separately; it appears in footnote~1 of Nisan and
Wigderson's communication-complexity paper \cite{NisanWigderson1995}.  An
explicit exposition is also given in \cite{HatamiKulkarniPankratov2011}.  Its
iterated composition has sensitivity \(6^k\) and degree \(3^k\).

This paper begins from that exact historical object.  Our first construction
has the desired fourteen sensitive directions but degree six rather than five.
Thus ``near construction'' has only the literal meaning ``one degree above the
target.''  It does not itself improve any exponent.  A degree-five,
sensitivity-fourteen function would clear the Kushilevitz exponent because
\[
  3^{18}=387420489>362797056=6^{11},\qquad
  14^{11}=4049565169664>3814697265625=5^{18}.
\]
The function constructed here has parameters \((6,14)\), not \((5,14)\).

Our results are structural and confined to the fixed constructions below.  We
identify the complete degree-six layer of \(P\), analyze selector and
parity-twist obstructions, and parameterize every Boolean correction relative
to \(P\) in two exact coordinate systems.  These parameterizations are not
ansatzes.  The obstruction theorems exclude only the explicitly stated
subfamilies within them.

\paragraph{Main results.}
\begin{enumerate}[label=(\roman*)]
  \item An explicit Kushilevitz-based Boolean function \(P\) on fourteen
  variables has sensitivity fourteen, degree six, and exactly 54 degree-six
  monomials, all with coefficient \(-1\) (\cref{thm:P}).
  \item A related degree-five, sensitivity-twelve base \(A\) has no arbitrary
  Boolean selector-neighbour \(B\) on the same variables with \(B(0)=1\) and
  \(\deg(B-A)\le4\) (\cref{thm:selector}).
  \item No target-labelled degree-five repair can factor only through the four
  block outputs and two ports (\cref{thm:gate-factor}).
  \item XOR twisting a gate-factored Boolean function by all four raw block
  parities always has degree greater than five
  (\cref{thm:parity-twist}).
  \item Additive and XOR correction systems are exact equivalences
  (\cref{thm:additive,thm:xor}).  For the natural representative \(J_0\), the
  corrected function \(P\oplus J_0\) has degree ten and an exact 81-term top
  layer (\cref{thm:J0}).
\end{enumerate}

\paragraph{Provenance of proof and computation.}
The construction formulas, noncancellation lemmas, selector obstruction,
gate-factor obstruction, parity-twist obstruction, and exact reformulations
are proved symbolically.  Exact finite computation is used only to materialize
truth tables and coefficient tables and to verify the complete coefficient
profile of \(P\oplus J_0\).  The top degree-ten formula also has a hand proof.
The accompanying verifier uses only integer arithmetic.

\section{Prior work and scope}\label{sec:related}

The distinction between the historical base function and the results of this
paper is important.

\paragraph{Kushilevitz and Nisan--Wigderson.}
The six-variable degree-three, sensitivity-six outer function is due to Eyal
Kushilevitz and was recorded by Nisan and Wigderson
\cite{NisanWigderson1995}.  Neither that function nor its recursive
amplification is claimed here as new.  Our displayed outer cubic is an explicit
coordinate relabeling of it.  Nisan and Wigderson identify their source as
``E.~Kushilevitz, private communication, 1994.''  What is studied here is a different, nonrecursive
substitution with four quadratic blocks and two linear ports, together with its
repair equations and obstruction theory.

\paragraph{Amano.}
Amano proved that no fully sensitive Boolean function on ten variables has
real degree at most four \cite[Theorem~3]{Amano2010}.  Restricting a function
at a point with at least ten sensitive coordinates to ten such coordinates
therefore gives the global consequence
\[
 \deg_{\mathbb R}(f)\le4\quad\Longrightarrow\quad s(f)\le9.
\]
Neither this bound nor Amano's ten-variable exclusion is claimed here as new.
Amano also reported a computer-verified obstruction for the narrower class of
fourteen-variable degree-five functions whose odd-degree coefficients lie in
\(\{0,1\}\) and even-degree coefficients lie in \(\{0,-1\}\)
\cite[Fact~6]{Amano2010}.  That bounded-polynomial obstruction and our
fixed-base obstructions are differently scoped; neither is asserted here to
subsume the other.  In later work, Amano enumerated, up to NPN equivalence,
nondegenerate Boolean functions of sensitivity three and determined extremal
values of several measures in that class \cite{Amano2017}.  Our
quadratic-rigidity proposition is only a local analytic statement under the
additional real-degree-two hypothesis, not a new enumeration or a
classification of all sensitivity-three functions.

\paragraph{Proskurin.}
Proskurin sharpened universal degree--block-sensitivity inequalities and
derived the unique univariate polynomial that could arise by symmetrizing a
hypothetical fully sensitive ten-variable degree-four function
\cite{Proskurin2021}.  Amano's earlier theorem already excludes the existence
of such a Boolean function.  Our claims neither improve Proskurin's universal
inequalities nor provide a new global low-degree sensitivity bound; they
concern specified twelve- and fourteen-variable Kushilevitz-based
constructions.

\paragraph{Maitra--Mukherjee--St\u{a}nic\u{a}--Tang.}
Maitra, Mukherjee, St\u{a}nic\u{a}, and Tang connected real degree to
resiliency, studied recursive amplification and higher-order sensitivity, and
analyzed low-degree fully sensitive functions on small numbers of variables
\cite{MaitraMukherjeeStanicaTang2025}.  Their computations include searches
among rotation-symmetric functions through ten variables.  Our
fourteen-variable construction does not improve their or Kushilevitz's
separation exponent.  The principal statements established in this manuscript
include the explicit 54-term degree-six component, the two exact fixed-base
parameterizations, and the selector, gate-factor, fixed-routing, parity-twist,
and local alternating-sum obstructions with their stated hypotheses.  A
sufficiently broad priority search has not yet been completed, so we make no
novelty or priority claim for these restricted statements.

\section{Definitions and the Kushilevitz outer function}

For \(x\in\cube n\) and \(i\in[n]\), write \(x^{\oplus i}\) for \(x\) with
coordinate \(i\) flipped.  The local and maximum sensitivities are
\[
 s(f,x)=\bigl|\{i:f(x)\ne f(x^{\oplus i})\}\bigr|,
 \qquad s(f)=\max_x s(f,x).
\]
All unqualified degrees below are real multilinear degrees.  The algebraic
normal form \(\ANF(f)\) is the unique multilinear polynomial over \(\Ftwo\);
its degree will always be written \(\deg_{\Ftwo}\), never simply ``degree.''
See \cite{ODonnell2014} for standard Boolean-function analysis conventions.

In zero-based coordinates, one standard form of the Kushilevitz cubic is
\[
 h_{\mathrm K}(z)
 =\sum_{i=0}^{5}z_i-\sum_{0\le i<j<6}z_iz_j
 +\sum_{D\in\mathcal D_{\mathrm K}}z_D,
\]
where \(z_D=\prod_{i\in D}z_i\) and
\[
\mathcal D_{\mathrm K}
=\{023,014,034,123,124,015,025,135,245,345\}.
\]
This is the function attributed to Kushilevitz in
\cite{NisanWigderson1995,HatamiKulkarniPankratov2011}.

For the block placement used in this paper it is convenient to relabel its
coordinates.  Put
\[
 K(y_0,\ldots,y_5)=h_{\mathrm K}(y_0,y_1,y_2,y_5,y_3,y_4).
\]
Equivalently,
\[
 K(y)=\sum_i y_i-\sum_{i<j}y_iy_j+\sum_{D\in\mathcal D}y_D,
\]
with
\[
 \mathcal D=\{013,014,024,025,035,123,125,145,234,345\}.
\]
Equivalently,
\((y_0,y_1,y_2,y_3,y_4,y_5)=(z_0,z_1,z_2,z_4,z_5,z_3)\).

\begin{proposition}[Kushilevitz function]\label{prop:K}
The set \(\mathcal D\) is a \(2\!-(6,3,2)\) design.  The function \(K\) is
Boolean, has degree exactly three, and satisfies
\(K(0)=0\) and \(K(e_i)=1\) for all six coordinates.
\end{proposition}

\begin{proof}
Direct inspection of the ten displayed triples shows that every pair occurs
exactly twice.  Hence every point occurs in five blocks.  If
\(S\subseteq[6]\), \(|S|=w\), and \(b(S)\) is the number of design blocks
contained in \(S\), then
\[
 K(1_S)=w-\binom w2+b(S).
\]
For \(w=0,1,2\) the values are \(0,1,1\).  For \(w=3\), the value is the
indicator that \(S\) is a block.  If \(w=4\), the complementary pair is met by
\(5+5-2=8\) blocks, leaving two blocks inside \(S\); the value is
\(4-6+2=0\).  If \(w=5\), exactly five blocks avoid the omitted point, so the
value is \(5-10+5=0\).  Finally, at \(w=6\) it is \(6-15+10=1\).
Thus \(K\) is Boolean.  Its ten cubic coefficients are nonzero, so its degree
is three, and the statements at the origin and singleton points are immediate.
\end{proof}

\section{Two block-signature noncancellation lemmas}\label{sec:signatures}

The obstruction proofs use two different direct-sum facts.  We state both in
full because conflating them conceals essential hypotheses.
We adopt the conventions \(\deg 0=-\infty\) and
\(\max\varnothing=-\infty\).

Let \(X_i=\{x_{i0},x_{i1},x_{i2}\}\) be pairwise disjoint triples and define
\[
 L_i=x_{i0}+x_{i1}+x_{i2},\qquad
 E_i=x_{i0}x_{i1}+x_{i0}x_{i2}+x_{i1}x_{i2},\qquad
 g_i=L_i-E_i=\NAE_3(X_i).
\]

\begin{lemma}[Weighted top signatures]\label{lem:weighted-signature}
Let \(z_1,\ldots,z_s\) be distinct raw coordinates disjoint from all blocks,
and let
\[
 Q(y,z)=\sum_{A\subseteq[r]}\sum_{B\subseteq[s]}
 c_{A,B}y_Az_B
\]
be the unique multilinear meta polynomial over a characteristic-zero field.
Give each \(y_i\) weight two and each \(z_j\) weight one.  For every
\(A,B\), and every choice of a pair \(p_i\in\binom{X_i}{2}\) for \(i\in A\),
the coefficient in \(Q(g,z)\) of
\[
 x_{\cup_{i\in A}p_i}z_B
\]
is exactly \((-1)^{|A|}c_{A,B}\).  Consequently,
\[
 \deg Q(g,z)=\max_{c_{A,B}\ne0}\bigl(2|A|+|B|\bigr).
\]
\end{lemma}

\begin{proof}
The displayed raw monomial contains exactly two variables from each block in
\(A\), none from any block outside \(A\), and exactly the ports in \(B\).  From
the meta monomial \(y_Az_B\), it is obtained only by choosing one quadratic
term from each \(-E_i\), and hence has coefficient
\((-1)^{|A|}c_{A,B}\).

No different meta support can contribute to the same monomial.  A block in
\(A'\setminus A\) would contribute at least one raw variable; a block in
\(A\setminus A'\) would contribute none; and a different port set has a
different raw port support.  Thus these top signatures lie in pairwise disjoint
monomial subspaces and cannot cancel.  A nonzero coefficient of maximum
weighted degree therefore produces a nonzero raw monomial of that degree,
while no substitution can exceed that weighted degree.  This proves the
formula.
\end{proof}

\begin{remark}
Disjointness and prior multilinearization are necessary.  If two formal gates
use the same raw block, \(y_0-y_1\) may vanish after substitution.  If a meta
polynomial is not multilinearized, \(y_0^2-y_0\) is an immediate counterexample.
The lemma also says nothing about a function that sees raw orientation inside a
block rather than only \(g_i\).
\end{remark}

For the second lemma use spin variables \(s_{ij}=1-2x_{ij}\), and put
\[
 \chi_i=s_{i0}s_{i1}s_{i2},\qquad
 u_i=\frac{s_{i0}s_{i1}+s_{i0}s_{i2}+s_{i1}s_{i2}}2,
 \qquad r_i=1-2g_i=u_i-\frac12.
\]

\begin{lemma}[Walsh block signatures]\label{lem:walsh-signature}
Let \(t_1,\ldots,t_s\) be port spins disjoint from the four blocks, and write a
multilinear meta-spin polynomial, after the shift \(r_i=u_i-\tfrac12\), as
\[
 q(r,t)=\sum_{U\subseteq[4]}\sum_{B\subseteq[s]}d_{U,B}u_Ut_B.
\]
After multiplication by \(\chi=\chi_0\chi_1\chi_2\chi_3\), every nonzero
\(d_{U,B}\) produces \(3^{|U|}\) distinct Walsh characters of degree
\[
 12-2|U|+|B|,
\]
each with coefficient \(2^{-|U|}d_{U,B}\).  Characters produced by distinct
\((U,B)\), or by different coordinate choices inside a block, are distinct.
Hence
\[
 \deg(\chi q)=\max_{d_{U,B}\ne0}\bigl(12-2|U|+|B|\bigr).
\]
\end{lemma}

\begin{proof}
For \(i\in U\),
\[
 \chi_i u_i=\frac{s_{i0}+s_{i1}+s_{i2}}2,
\]
so block \(i\) contributes one of three degree-one characters.  For
\(i\notin U\), it contributes the unique degree-three character \(\chi_i\).
The ports contribute exactly the character indexed by \(B\).

The blockwise signature ``one selected spin'' versus ``all three spins''
recovers \(U\), the selected spin recovers the coordinate choice, and the port
support recovers \(B\).  The corresponding Walsh characters are therefore
pairwise distinct.  Since Walsh characters form a basis, no cross-signature
cancellation is possible.  Summing the four block degrees and the port degree
gives the stated degree formula.
\end{proof}

\section{The degree-six, sensitivity-fourteen composition}

Use four disjoint blocks
\(X_i=(x_{3i},x_{3i+1},x_{3i+2})\), \(0\le i<4\), and define
\[
 P=K(g_0,g_1,g_2,g_3,x_{12},x_{13}).
\]

\begin{theorem}\label{thm:P}
The function \(P\) is Boolean, satisfies
\[
 P(0)=0,\qquad P(e_j)=1\quad(0\le j<14),
\]
and therefore has sensitivity fourteen at the origin.  Its degree is exactly
six, and its homogeneous degree-six part is
\[
 P_6=-E_0E_1E_3-E_1E_2E_3=-E_1E_3(E_0+E_2).
\]
It contains exactly 54 degree-six monomials, each with coefficient \(-1\), and
has no term of higher degree.
\end{theorem}

\begin{proof}
Booleanity follows by composing Boolean functions.  At the raw origin all six
outer inputs vanish.  A singleton in one of the first twelve coordinates
changes exactly one \(g_i\) from zero to one; the final two coordinates are
the two remaining outer singleton directions.  The target rows follow from
\cref{prop:K}.

By \cref{lem:weighted-signature}, raw degree six can arise only from an outer
cubic containing three heavy inputs.  In \(\mathcal D\) the only such triples
are \(013\) and \(123\), whose top terms are respectively
\(-E_0E_1E_3\) and \(-E_1E_2E_3\).  Their block signatures differ, so they
cannot cancel.  Each product of three \(E_i\)'s has \(3^3=27\) distinct
monomials; the two sets are disjoint.  No outer term has weighted degree above
six.
\end{proof}

\section{The central selector obstruction}\label{sec:selector}

The following twelve-variable slice is a natural attempt to obtain a
degree-five base and add one or two selector variables.  Use three NAE blocks
on \(x_0,\ldots,x_8\), put
\(\ell_3=x_9,\ell_4=x_{10},\ell_5=x_{11}\), and define
\[
 A=K(g_0,g_1,g_2,\ell_3,\ell_4,\ell_5).
\]

\begin{theorem}[Arbitrary-selector obstruction]\label{thm:selector}
The function \(A\) is Boolean, has sensitivity twelve at the origin, and has
degree exactly five.  Its degree-five part is
\[
 A_5=E_0E_1(\ell_3+\ell_4)
     +E_0E_2(\ell_4+\ell_5)
     +E_1E_2(\ell_5+\ell_3),
\]
which consists of 54 distinct coefficient-\(+1\) monomials.

There is no Boolean function \(B:\cube{12}\to\{0,1\}\) satisfying
\[
 B(0)=1,
 \qquad \deg(B-A)\le4.
\]
\end{theorem}

\begin{proof}
The heavy triple \(012\) is not a block of \(\mathcal D\), so no outer cubic
contains three heavy inputs.  The six heavy-heavy-linear blocks are
\[
 013,014,024,025,125,123.
\]
The weighted-signature lemma yields the displayed degree-five part.  Each of
its six summands contributes nine distinct masks, distinguished by their heavy
block pair and linear port.  The Booleanity and twelve singleton rows follow as
in \cref{thm:P}.

Assume that \(B\) exists.  Average \(B\) under independent coordinate
permutations inside each heavy block and under independently complementing all
three coordinates of each block.  Call the average \(C\).  These affine cube
automorphisms preserve real degree and fix \(A\); therefore
\[
 0\le C\le1,\qquad \deg(C-A)\le4.
\]
The action on one raw block has exactly two orbits, the points with NAE value
zero and those with NAE value one.  Thus \(C\) factors through the three NAE
outputs and three linears.  The meta-origin orbit contains the eight points
obtained by choosing each raw block to be \(000\) or \(111\).  Since it contains
the raw origin and \(B(0)=1\),
\[
 c_0=C(0,0,0,0,0,0)\ge\frac18.
\]

Apply \cref{lem:weighted-signature} to the unique multilinear meta polynomial
of \(C-A\).  Every meta coefficient of weight greater than four vanishes.
Consequently all HHH coefficients of \(C\), including their linear supersets,
are zero; the six HHL coefficients agree with those of \(A\) and equal one;
and all HHLL and HHLLL coefficients vanish.

Fix a linear assignment \(z\in\cube3\), and write \(h_z\) for \(C\) as a
function of the three heavy variables.  It is quadratic.  If
\(b_{ij}(z)\) are its three heavy-pair coefficients, the preceding coefficient
identities give
\[
\begin{aligned}
 b_{01}(z)&=\beta_{01}+z_3+z_4,\\
 b_{02}(z)&=\beta_{02}+z_4+z_5,\\
 b_{12}(z)&=\beta_{12}+z_5+z_3.
\end{aligned}
\]
For a quadratic polynomial on three variables define
\[
 S_z=h_z(111)-h_z(100)-h_z(010)-h_z(001)+2h_z(000).
\]
M\"obius inversion on the three-cube shows that \(S_z\) is the sum of the
three pair coefficients.  Hence
\[
 S_{111}=S_{000}+6.
\]
Because every value of \(h_z\) lies in \([0,1]\),
\[
 S_{000}\ge-3+2c_0,
 \qquad S_{111}\le3.
\]
But \(c_0\ge1/8\) implies
\[
 S_{111}\ge-3+\frac14+6=\frac{13}{4}>3,
\]
a contradiction.
\end{proof}

\begin{remark}[Exact scope]
The theorem treats arbitrary Boolean \(B\) on the same twelve raw variables;
it is not a restriction to gate-factored selectors.  It depends essentially on
the fixed base \(A\), the anchored value \(B(0)=1\), and the degree-four
difference.  For example, without \(B(0)=1\) one may take \(B=A\); without
Booleanity one may take \(B=A+1\).  It does not exclude a different base,
additional variables, or a different outer architecture.
\end{remark}

\section{Gate-factor and parity-twist obstructions}\label{sec:central-obstructions}

We next return to the fourteen-variable setting.  These results explain why a
successful repair cannot see only the four NAE output bits and why the most
obvious full-block parity dualization also fails.

\begin{proposition}[Quadratic rigidity]\label{prop:quadratic}
Every Boolean function of real degree at most two has sensitivity at most
three.  More precisely, if \(h(0)=0\) and its sensitive set at the origin is
exactly \(\{a,b,c\}\), then
\[
 h=\NAE_3(x_a,x_b,x_c)
\]
and \(h\) is independent of every other coordinate.
\end{proposition}

\begin{proof}
Translate a maximally sensitive point to the origin and complement the output
if necessary; these operations preserve degree.  Write
\[
 h=\sum_i\alpha_i x_i+\sum_{i<j}\beta_{ij}x_ix_j.
\]
Every sensitive coordinate has \(\alpha_i=1\).  On a pair of sensitive
coordinates, Booleanity gives \(2+\beta_{ij}\in\{0,1\}\), so
\(\beta_{ij}\in\{-2,-1\}\).  On any sensitive triple, a coefficient \(-2\)
would force the triple value to be negative; hence all three pair coefficients
are \(-1\).  Four sensitive coordinates are then impossible, because their
all-one face value would be \(4-\binom42=-2\).

In the equality case, the restriction to \(a,b,c\) is NAE.  For an outside
coordinate \(k\), each \(\beta_{ik}\) with \(i\in\{a,b,c\}\) belongs to
\(\{-1,0\}\).  Evaluating at \(e_a+e_b+e_c+e_k\) forces all three to zero.
For two outside coordinates \(k,\ell\), Booleanity gives
\(\beta_{k\ell}\in\{0,1\}\), while evaluation at
\(e_a+e_k+e_\ell\) forces it to zero.
\end{proof}

\begin{theorem}[Gate-factor obstruction]\label{thm:gate-factor}
Let \(T:\cube6\to\{0,1\}\) be arbitrary and put
\[
 F=T(g_0,g_1,g_2,g_3,x_{12},x_{13}),
\]
where the four raw blocks and the two coordinate ports are pairwise disjoint.
If \(F(0)=0\) and \(F(e_j)=1\) for every \(0\le j<14\), then
\(\deg F>5\).
\end{theorem}

\begin{proof}
All 64 meta inputs are independently attainable, so \(T\) has a unique
multilinear meta representation.  If \(\deg F\le5\),
\cref{lem:weighted-signature} forces every meta coefficient of weight above
five to vanish.  Set both ports to zero.  The remaining Boolean function of
\(g_0,g_1,g_2,g_3\) has meta degree at most two, value zero at its origin, and
value one at all four meta singletons.  This contradicts
\cref{prop:quadratic}.
\end{proof}

\begin{remark}
The target labels are essential: the constant function is a gate-factored
counterexample without them.  The theorem does not cover raw orientation
dependence inside a block.
\end{remark}

Let \(p_i=x_{i0}\oplus x_{i1}\oplus x_{i2}\) denote XOR parity on raw block
\(i\).

\begin{theorem}[Parity twist]\label{thm:parity-twist}
For every Boolean meta function
\(T(g_0,g_1,g_2,g_3,x_{12},x_{13})\), the Boolean function
\[
 F=T\oplus p_0\oplus p_1\oplus p_2\oplus p_3
\]
has real degree greater than five.
\end{theorem}

\begin{proof}
Passing between bit variables and spin variables preserves the degree
filtration.  Let \(q_T=1-2T\).  Expand \(q_T\) uniquely in the shifted block
variables \(r_i=u_i-\tfrac12\) and the two port spins.  The output spin of
\(F\) is
\[
 q_F=q_T\chi_0\chi_1\chi_2\chi_3.
\]
If \(\deg q_F\le5\), \cref{lem:walsh-signature} forces every sector with
\(|U|\le3\) to vanish, since its core degree is at least six.  In the only
remaining sector \(U=[4]\), the coefficient has port degree at most one.
Thus
\[
 q_T=a(t_{12},t_{13})u_0u_1u_2u_3
\]
for an affine function \(a\) of the port spins.

Fix the ports.  If \(a=0\), then \(q_T=0\), already contradicting Booleanity.
Otherwise the right side has magnitude \(|a|\prod_i|u_i|\).  For one block,
\(|u_i|=3/2\) on an all-equal raw
orientation and \(|u_i|=1/2\) on a non-all-equal orientation.  All sixteen
meta patterns are attainable, so the magnitude varies with the block outputs.
This contradicts Booleanity, which requires \(|q_T|=1\) everywhere.
\end{proof}

\begin{remark}[Exact scope]
This theorem concerns a meta function followed by the parity twist on all four
blocks.  It proves nothing about a partial twist, a nonlinear orientation
correction, or a \(T\) that already sees raw orientations.  Indeed, if raw
orientation is allowed, choosing \(T=p_0\oplus p_1\oplus p_2\oplus p_3\)
gives the constant zero function after twisting.
\end{remark}

\section{Two exact reformulations of every repair}\label{sec:exact-systems}

The preceding theorems exclude particular structural classes.  We now record
two descriptions that exclude nothing.  They are changes of coordinates on
the full set of Boolean functions relative to the fixed function \(P\), not
search restrictions.

Let
\[
 \mathcal T=\{0,e_0,e_1,\ldots,e_{13}\}
\]
be the fifteen target rows, and put
\[
 \Delta=E_1E_3(E_0+E_2)=-P_6.
\]

\begin{theorem}[Exact additive residual system]\label{thm:additive}
Let \(R:\cube{14}\to\mathbb R\), let the same symbol denote its unique real
multilinear polynomial, and set \(F=P+R\).  Then \(F\) is Boolean, has degree
at most five, and agrees with \(P\) on \(\mathcal T\) if and only if all three
of the following hold:
\begin{align}
 R&=\Delta+L &&\text{for a multilinear polynomial }\deg L\le5,
                                                        \label{eq:add-high}\\
 R(x)\bigl(R(x)+2P(x)-1\bigr)&=0
       &&\text{for every }x\in\cube{14},                 \label{eq:add-bool}\\
 R(t)&=0 &&\text{for every }t\in\mathcal T.              \label{eq:add-label}
\end{align}
Consequently, every admissible Boolean \(F\) occurs exactly once, with
\(R=F-P\).
\end{theorem}

\begin{proof}
Since \(P\) has no terms above degree six and \(P_6=-\Delta\), cancellation
of every term of degree at least six in \(P+R\) is equivalent to
\cref{eq:add-high}.  This proves the degree equivalence, including the absence
of any term of degree greater than six in \(R\).

On a cube point, \(P^2=P\).  Therefore
\[
 F(F-1)=(P+R)(P+R-1)=R(R+2P-1).
\]
The left side vanishes exactly when \(F\in\{0,1\}\), proving
\cref{eq:add-bool}.  More explicitly, where \(P=0\), the permitted residual
values are \(R\in\{0,1\}\), and where \(P=1\), they are
\(R\in\{0,-1\}\).  Finally, \(F(t)=P(t)\) is exactly \(R(t)=0\), which proves
\cref{eq:add-label}.  The identity \(R=F-P\) proves uniqueness.
\end{proof}

All three conditions are indispensable.  The residuals \(R=-P\) and
\(R=1-P\) satisfy the Boolean equation but give the constant-zero and
constant-one functions, so they demonstrate the need for the target-row
condition.  Conversely, \(R=\Delta\) satisfies the high-degree cancellation
and vanishes on \(\mathcal T\), but at
\(x=e_0+e_1+e_3+e_4+e_9+e_{10}\) one has
\(P(x)=\Delta(x)=1\), hence \((P+R)(x)=2\).  This demonstrates the need for
the Boolean equation.

\begin{theorem}[Exact XOR correction system]\label{thm:xor}
For every Boolean \(F:\cube{14}\to\{0,1\}\) there is a unique Boolean mask
\[
 J=P\oplus F,
\]
and pointwise, as an identity over the reals,
\begin{equation}
 F=P+(1-2P)J.                                             \label{eq:xor-real}
\end{equation}
Thus the Boolean functions of degree at most five that agree with \(P\) on
\(\mathcal T\) are in bijection with the Boolean masks \(J\) satisfying
\begin{equation}
 J|_{\mathcal T}=0,\qquad
 \deg_{\mathbb R}\bigl(P+(1-2P)J\bigr)\le5.              \label{eq:xor-exact}
\end{equation}

Every such mask necessarily satisfies
\begin{equation}
 \ANF_{\ge6}(J)=\ANF_{\ge6}(P),                           \label{eq:anf-high}
\end{equation}
which here consists of exactly the 54 degree-six masks in
\(E_1E_3(E_0+E_2)\), and no higher mask.
\end{theorem}

\begin{proof}
The first assertion is the pointwise identity
\(a\oplus b=a+(1-2a)b\) for \(a,b\in\{0,1\}\).  It also gives
\(J=P\oplus F\), so the representation is unique.  Equality of the target
rows is equivalent to \(J=0\) there, proving the exact bijection
\cref{eq:xor-exact}.

For a Boolean function, reducing every integer coefficient of its real
M\"obius expansion modulo two gives its \(\Ftwo\)-algebraic normal form.  If
\(\deg_{\mathbb R}F\le5\), all real coefficients of \(F\) on masks of size at
least six vanish and therefore so do the corresponding ANF coefficients.
Since \(F=P\oplus J\), its ANF is the sum of the ANFs of \(P\) and \(J\).
This proves \cref{eq:anf-high}.  The description of the high masks follows
from \cref{thm:P} after reduction of \(-1\) modulo two.
\end{proof}

Equation~\eqref{eq:anf-high} immediately shows that any successful \(J\)
depends on raw variables from all four NAE blocks.  It is nevertheless only a
necessary parity condition; real-degree cancellation is stricter.

Define the compact Boolean mask
\begin{equation}
 J_0=g_1g_3(g_0\oplus g_2).                               \label{eq:J0}
\end{equation}
Its degree-six ANF layer is precisely
\(E_1E_3(E_0+E_2)\), and it has no higher ANF term.  Hence every mask
satisfying \cref{eq:anf-high} can be written uniquely as
\[
 J=J_0\oplus W,
 \qquad \deg_{\Ftwo}W\le5.
\]
This is again a reparameterization, not an assertion that \(W\) is local,
factored, or otherwise simple.  Theorems~\ref{thm:gate-factor} and
\ref{thm:parity-twist} explain why two especially natural simplifications do
not suffice.

\section{The exact failure of the simplest mask \texorpdfstring{\(J_0\)}{J0}}
\label{sec:J0}

The degree of \(J_0\) must not be confused with the degree of the repaired
function.  The former is eight; the latter is ten.

\begin{theorem}[The function \(P\oplus J_0\) has degree ten]\label{thm:J0}
The real multilinear degree of \(J_0\) is eight, with top component
\[
 (J_0)_8=-2E_0E_1E_2E_3.
\]
For
\[
 F_0=P\oplus J_0
\]
the real multilinear degree is ten, and
\[
 (F_0)_{10}=4E_0E_1E_2E_3x_{12}x_{13}.
\]
Thus \(F_0\) has exactly \(3^4=81\) degree-ten monomials, each with
coefficient \(+4\).
\end{theorem}

\begin{proof}
At the six meta inputs \(y=(y_0,y_1,y_2,y_3,y_4,y_5)\), the real polynomial
of \(J_0\) is
\[
 j=y_1y_3(y_0+y_2-2y_0y_2).
\]
The unique weight-eight term after substituting the four heavy NAE gates is
\(-2g_0g_1g_2g_3\), whose top component is
\(-2E_0E_1E_2E_3\).  Weighted top-signature noncancellation proves that it is
nonzero and that no term has larger degree.

Multilinearizing the meta function \(K\oplus j=K+j-2Kj\) gives
\begin{align*}
 K
 &-y_0y_1y_3-y_1y_2y_3+2y_0y_1y_2y_3\\
 &+2y_0y_1y_3y_4+2y_1y_2y_3y_4
  +2y_0y_1y_3y_5+2y_1y_2y_3y_5\\
 &-4y_0y_1y_2y_3y_4-4y_0y_1y_2y_3y_5\\
 &-2y_0y_1y_3y_4y_5-2y_1y_2y_3y_4y_5
  +4y_0y_1y_2y_3y_4y_5.
\end{align*}
Only the final meta monomial has weighted degree ten.  Substituting
\(y_i=g_i\) for \(0\le i<4\), \(y_4=x_{12}\), and \(y_5=x_{13}\), its top
signature is
\(4E_0E_1E_2E_3x_{12}x_{13}\).  The weighted-signature lemma proves both
noncancellation and the absence of any higher term.  Each \(E_i\) offers three
pair choices on a disjoint block, giving exactly \(3^4=81\) monomials.
\end{proof}

For completeness, exact integer M\"obius inversion gives the following full
nonzero coefficient profile.  An entry \(c:m\) means that exactly \(m\)
monomials of that degree have coefficient \(c\).
\begin{center}
\begin{tabular}{ccl}
\toprule
degree & number nonzero & coefficient multiplicities\\
\midrule
1  & 14   & \(+1:14\)\\
2  & 91   & \(-1:91\)\\
3  & 192  & \(+1:192\)\\
4  & 357  & \(-1:168,\ +2:189\)\\
5  & 918  & \(-4:162,\ -2:702,\ +1:54\)\\
6  & 1701 & \(+2:972,\ +4:729\)\\
7  & 1890 & \(-2:594,\ -4:1296\)\\
8  & 1269 & \(+2:135,\ +4:1134\)\\
9  & 486  & \(-4:486\)\\
10 & 81   & \(+4:81\)\\
\bottomrule
\end{tabular}
\end{center}
The total is 6999 nonzero coefficients.  The table is a reproducible finite
calculation, whereas the degree-ten and 81-term assertions in
\cref{thm:J0} are proved independently of it.

\section{Further exact obstructions and residual constraints}

The central no-go theorems already show that a successful correction, if one
exists, cannot be a selector extension of \(A\), a target-labelled function of
the six meta outputs alone, or a full four-block parity twist of such a meta
function.  Two additional consequences clarify the remaining scope.

\begin{corollary}[Rigidity of the fixed outer routing]
For \(0\le i<4\), let \(h_i\) be a Boolean function of real degree at most two
that vanishes at its own origin and is sensitive there in all three
coordinates of a designated raw-coordinate triple \(X_i\).  Assume the
\(X_i\) are pairwise disjoint and route \(h_i\) to input \(i\) of \(K\).
Route two further raw coordinates, distinct from one another and from every
\(X_i\), to inputs \(4\) and \(5\).  Then, after relabeling within blocks, the
four quadratic gates are exactly the four NAE gates used in \(P\).  The
composite therefore has the nonzero degree-six layer in \cref{thm:P} and is
sensitive in all fourteen designated directions.
\end{corollary}

\begin{proof}
Each quadratic gate must be sensitive in all three of its designated
coordinates.  Proposition~\ref{prop:quadratic} makes it NAE on exactly that
triple and independent of every other coordinate.  The claimed composite is
therefore \(P\) up to within-block relabeling, which preserves its degree-six
top signatures.
\end{proof}

There is also a local constraint on every exact correction.  For either outer
design block \(013\) or \(123\), choose one of the three quadratic monomials
from each associated \(E_i\).  Fix all other raw variables, including both
ports, to zero.  On the resulting six-dimensional face, each restricted NAE
gate is \(\OR_2\), and the restriction of \(P\) is \(\OR_6\).  There are 27
faces of each block signature, hence 54 such faces.

\begin{proposition}[Alternating zero-set constraint]\label{prop:or6}
Let \(F:\cube6\to\{0,1\}\) have degree at most five, with \(F(0)=0\), and
let
\[
 Z=\{x\ne0:F(x)=0\}.
\]
Then
\[
 \sum_{x\in Z}(-1)^{|x|}=-1.
\]
Consequently, every degree-five target-preserving correction of \(P\) must
satisfy this identity on each of the 54 \(\OR_6\) faces above.
\end{proposition}

\begin{proof}
The coefficient of the full six-variable monomial is
\[
 \sum_{x\in\cube6}(-1)^{6-|x|}F(x)
 =\sum_{x\in\cube6}(-1)^{|x|}F(x).
\]
It vanishes because \(\deg F\le5\).  Also
\(\sum_{x\in\cube6}(-1)^{|x|}=0\), while the origin is not in \(Z\).
Subtracting the zero-set contribution from the full alternating sum proves
the formula.  Restriction cannot increase degree, so it applies to every
listed face.
\end{proof}

This condition is necessary, not sufficient.  In particular, it does not
replace either exact system in \cref{sec:exact-systems}.

\section{Standalone exact verifier}\label{sec:verifier}

The supplementary script \texttt{verify\_exact.py} is standalone, uses only
the Python standard library, and has no numerical or symbolic-algebra
dependency.  It constructs \(K,P,A,J_0\), and \(P\oplus J_0\) from the
displayed formulas.  For a truth vector \(v[0],\ldots,v[2^n-1]\) in numeric
mask order, it performs the in-place integer subset transform
\[
 c\leftarrow v;\qquad
 c[m]\leftarrow c[m]-c[m\setminus\{i\}]
 \quad(i=0,\ldots,n-1;\ i\in m).
\]
The resulting \(c[m]\) are exactly the real multilinear M\"obius
coefficients; no division or floating-point operation occurs.

The verifier computes every real multilinear coefficient and checks the
claimed Booleanity, origin sensitivities, degrees, top supports, and the
complete coefficient histogram of \(P\oplus J_0\).  Its deterministic summary
is:
\begin{center}
\begingroup
\small
\setlength{\tabcolsep}{5pt}
\begin{tabular}{lrrrrl}
\toprule
function & variables & degree & origin sensitivity & nonzero coefficients
& top layer\\
\midrule
\(K\)            & 6  & 3  & 6  & 31   & ten coefficient-\(+1\) cubics\\
\(A\)            & 12 & 5  & 12 & 421  & 54 coefficients \(+1\)\\
\(P\)            & 14 & 6  & 14 & 951  & 54 coefficients \(-1\)\\
\(J_0\)          & 14 & 8  & 0  & 1728 & 81 coefficients \(-2\)\\
\(P\oplus J_0\) & 14 & 10 & 14 & 6999 & 81 coefficients \(+4\)\\
\bottomrule
\end{tabular}
\endgroup
\end{center}
The script exits nonzero upon failure of any encoded assertion.  The
diagnostic row for \(J_0\) is included specifically to prevent the erroneous statement
\(\deg J_0=10\): the degree-ten object is \(P\oplus J_0\).

\section{What is and is not resolved}

The two parameterizations are exact and complete relative to the fixed base
\(P\), but the obstruction theorems exclude only their stated subfamilies.
They leave open whether there exists a Boolean mask \(J\), equivalently an
additive residual \(R\), satisfying the systems in
\cref{thm:additive,thm:xor}.  Any such correction must obey the 54 forced ANF
masks, the target rows, the real Boolean equations, the degree cancellation,
and the local alternating constraints.  It must also evade the selector,
gate-factor, fixed-routing, and full-parity-twist theorems.  These facts point
to raw orientation-dependent interactions across blocks, but they do not
constitute an existence theorem, a completeness result for any smaller
family, or a lower bound for all degree-five functions.

In particular, this paper proves no degree-five, sensitivity-fourteen
construction and therefore no new asymptotic separation.  The phrase
``near-solution'' refers only to the explicit degree-six function \(P\) and
its explicitly computed degree-six homogeneous component.

\section*{Generative-AI disclosure}

Generative-AI systems were used extensively for conjecture generation,
symbolic exploration, code generation, proof drafting, internal auditing, and
manuscript editing.  The named human authors reviewed the final theorem
statements, proofs, computations, and references and accept responsibility for
the contents.

\bibliographystyle{alpha}
\bibliography{references}

\end{document}
